\chapter{Cradle Cap}

\section*{Definition and Overview}
Cradle cap (seborrhoeic dermatitis of infancy) is a common, harmless condition in which babies develop greasy, yellowish, scaly patches on their scalp. It usually appears in the first few months of life and typically clears on its own within weeks to months. Cradle cap can also affect the eyebrows, ears, face, and nappy area. It is not itchy or painful for most babies and is not caused by poor hygiene or an allergy. Treatment is usually simple: gentle washing, soft brushing, and, if needed, a mild emollient or medicated product.

\section*{Epidemiology}
Cradle cap is very common, affecting about 1 in 10 babies in the first months of life, with some estimates higher. It usually begins between two and six weeks of age and resolves by the time the baby is six to twelve months old. It affects boys and girls equally. A family history of eczema, asthma, or seborrhoeic dermatitis may increase the chance, but most babies with cradle cap are otherwise completely healthy.

\section*{Aetiology and Pathophysiology}
The exact cause is not fully understood, but it involves the activity of the skin's oil glands and a common yeast.

\begin{itemize}
    \item \textbf{Sebaceous glands:} babies' oil glands are active due to maternal hormones, increasing sebum production
    \item \textbf{Malassezia yeast:} a common skin yeast that thrives in oily skin and may contribute to flaking
    \item \textbf{Scale formation:} excess oil and skin cells clump into greasy scales
    \item \textbf{Not an allergy:} cradle cap is not caused by a food allergy or contact allergy
    \item \textbf{Not infectious:} it cannot be passed to other children or adults
    \item \textbf{Self-limiting:} it resolves as the hormone effects fade and the oil glands settle
\end{itemize}

\section*{Clinical Features}
Cradle cap has a characteristic appearance.

\begin{itemize}
    \item Greasy, yellowish, or whitish scales on the scalp
    \item Crusty patches that may flake off
    \item Mild redness under the scales in some babies
    \item Similar scaling on the eyebrows, behind the ears, in skin folds, or in the nappy area
    \item Generally no itching or discomfort
    \item Normal hair growth and general health
    \item Symptoms that come and go but usually resolve by the first birthday
\end{itemize}

\section*{Differential Diagnosis}
Most scaly scalps in babies are cradle cap, but other conditions are considered.

\begin{itemize}
    \item Atopic dermatitis (eczema), which is itchy and often affects the face and body
    \item A fungal infection (tinea capitis), which causes patches of hair loss
    \item Psoriasis, rare in infancy
    \item Impetigo or other bacterial infections, with crusting and weeping
    \item Nappy rash, when the rash spreads beyond the scalp
    \item Lice or other infestations, which are uncommon in babies
\end{itemize}

\section*{When to Seek Medical Care}
Cradle cap does not usually need medical treatment, but see a GP, child health nurse, or paediatrician if the rash spreads, becomes red and inflamed, weeps, or is associated with hair loss, fever, or an unwell baby.

\textbf{Call 000 (or local emergency services) for:}
\begin{itemize}
    \item The baby develops a high fever or becomes very drowsy, floppy, or difficult to wake
    \item The baby is struggling to breathe or appears seriously unwell
    \item A seizure or loss of consciousness
\end{itemize}

Seek urgent medical review if the rash becomes infected with weeping or spreading redness, there is significant hair loss, or the baby seems generally unwell.

\section*{Diagnosis and Investigations}
Diagnosis is clinical and needs no tests in typical cases.

\begin{itemize}
    \item \textbf{History and examination:} the characteristic appearance and the baby's age
    \item \textbf{Assessment of the whole skin:} to check for eczema or nappy rash
    \item \textbf{No routine tests:} blood tests and skin scrapings are not needed unless the diagnosis is unclear
    \item \textbf{Review:} a check that the rash resolves as expected
\end{itemize}

\section*{Management}
Treatment is simple and aimed at loosening and removing the scales.

\begin{itemize}
    \item \textbf{Gentle washing:} wash the scalp regularly with baby shampoo and warm water
    \item \textbf{Soft brushing:} brush the scalp with a soft baby brush to lift the scales
    \item \textbf{Oil application:} massage a little baby oil, olive oil, or emollient into the scalp, leave it for a short time, then wash it out
    \item \textbf{Emollients:} a mild moisturiser for dry, flaky skin
    \item \textbf{Medicated products:} a gentle anti-dandruff or ketoconazole shampoo, or a low-strength steroid cream, if advised by a doctor
    \item \textbf{Patience:} cradle cap often recurs and settles over time
\end{itemize}

\section*{Complications}
\begin{itemize}
    \item Secondary bacterial infection from scratching or aggressive scrubbing
    \item Skin irritation from over-treatment or harsh products
    \item Mild eczema developing in children who are prone to it
    \item Anxiety in parents, although the condition is harmless
    \item Persistent scaling beyond infancy in some children
\end{itemize}

\section*{Prognosis and Outcomes}
Cradle cap is a harmless, self-limiting condition. In the vast majority of babies it clears by six to twelve months of age, whether or not it is treated. Treatment speeds up the clearing of scales but does not prevent recurrence. There are no long-term consequences, although some children who had cradle cap develop seborrhoeic dermatitis or eczema later in childhood.

\section*{Prevention and Self-Care}
\begin{itemize}
    \item Wash the baby's scalp regularly with a gentle baby shampoo
    \item Brush the scalp with a soft brush to prevent scale buildup
    \item Use oils and emollients sparingly, and rinse thoroughly
    \item Avoid harsh scrubbing, picking, or pulling at the scales
    \item Keep the skin moisturised in other affected areas
    \item Seek advice if the rash worsens or looks infected
\end{itemize}

\section*{Sources and Further Reading}
The following reputable sources provide further information for healthcare students and parents seeking reliable, current advice about cradle cap.

\begin{itemize}
    \item Healthdirect Australia. \textit{Cradle cap.} https://www.healthdirect.gov.au/
    \item Raising Children Network. \textit{Cradle cap.} https://raisingchildren.net.au/
    \item NHS. \textit{Cradle cap.} https://www.nhs.uk/conditions/cradle-cap/
    \item Royal Children's Hospital Melbourne. \textit{Kids Health Info: cradle cap.} https://www.rch.org.au/kidsinfo/
    \item Mayo Clinic. \textit{Cradle cap: overview.} https://www.mayoclinic.org/diseases-conditions/cradle-cap/
    \item MedlinePlus. \textit{Cradle cap.} https://medlineplus.gov/ency/article/002078.htm
\end{itemize}
